\section{Análisis y conclusiones}

En esta sección, resuman brevemente lo realizado en el laboratorio. Indiquen qué se hizo y cuál fue el propósito de cada actividad, haciendo referencia a las secciones, tablas y ecuaciones correspondientes cuando sea necesario. No es necesario volver a enunciar todos los datos o repetir los procedimientos detallados anteriormente. El resumen puede realizarse en un solo párrafo o en tres párrafos distintos, uno para cada parte.

Luego del resumen de las actividades, respondan las preguntas de análisis planteadas en las tres partes de la guía. Pueden presentar todas las respuestas en un solo texto o separarlas en tres párrafos, uno para cada actividad. En cada respuesta, no se limiten a indicar un resultado: expliquen el razonamiento que utilizaron para llegar a él y relacionen sus conclusiones con los datos, cálculos y gráficos obtenidos durante el trabajo. Las respuestas a las preguntas de análisis deben estar integradas en el texto, conectando ideas mediante conectores lógicos. Eviten presentar las respuestas de manera aislada, como en formato pregunta-respuesta.

Finalmente, incluyan las principales conclusiones a las que llegaron como grupo. Estas deben sintetizar los aprendizajes obtenidos a partir de las actividades realizadas y destacar los aspectos que consideren más relevantes respecto de la medición de magnitudes físicas, el tratamiento de datos y su representación gráfica. Cuando corresponda, indiquen también las principales limitaciones o fuentes de incertidumbre que hayan identificado durante el trabajo.

Las conclusiones deben basarse en los resultados obtenidos por el grupo. \textbf{No incluyan citas bibliográficas en esta sección}.

\textbf{\underline{Extensión máxima del informe:} indefinida}.
